% Circuito ativo estudado
\begin{frame}{Filtro passa-faixa de realimentação múltipla}
  \begin{columns}[T,onlytextwidth]
    \begin{column}{0.60\textwidth}
      \centering
      \includegraphics[
        page=7,
        trim=135bp 572bp 120bp 78bp,
        clip,
        width=\linewidth,
        height=0.61\textheight,
        keepaspectratio
      ]{\detokenize{prism-uploads/Romário_Jonas_-_relatório_Pratica_02___Circuitos_Elétricos_2.pdf}}

      \vspace{0.08cm}
      \scriptsize Circuito MFB analisado.
    \end{column}
    \hfill
    \begin{column}{0.36\textwidth}
      \small
      \begin{block}{Topologia}
        Filtro ativo passa-faixa de segunda ordem, com amplificador operacional e
        duas malhas de realimentação.
      \end{block}

      \vspace{0.12cm}
      \begin{block}{Valores nominais}
        \begin{itemize}
          \item $R_1=R_2=470\,\Omega$;
          \item $R_3=47\,\mathrm{k}\Omega$;
          \item $C_1=C_2=100\,\mathrm{nF}$;
          \item alimentação: $\pm 10\,\mathrm{V}$;
          \item entrada: senoide de $100\,\mathrm{mV}$.
        \end{itemize}
      \end{block}
    \end{column}
  \end{columns}
\end{frame}